% CENG 467 — Natural Language Understanding and Generation, Spring 2026
% Term Project — Progress Report (Week 11 Technical Checkpoint)
%
% Build: pdflatex progress_report.tex  (no bibtex needed)

\documentclass[11pt,a4paper]{article}

\usepackage[margin=2.5cm]{geometry}
\usepackage[T1]{fontenc}
\usepackage{lmodern}
\usepackage{microtype}
\usepackage{titling}
\usepackage{enumitem}
\usepackage{booktabs}
\usepackage{xcolor}
\usepackage{hyperref}
\usepackage{listings}
\usepackage{amsmath,amssymb}

\hypersetup{
  colorlinks=true,
  linkcolor=black,
  urlcolor=blue!60!black,
  citecolor=blue!60!black,
}

\lstset{
  basicstyle=\ttfamily\small,
  breaklines=true,
  frame=single,
  framesep=4pt,
  xleftmargin=4pt,
  xrightmargin=4pt,
  backgroundcolor=\color{gray!8},
}

\title{\vspace{-1.5cm}\textbf{Progress Report} \\[0.3em]
\large Controllable Text Generation via Diffusion Models}
\author{Zübeyr Almaho (300201023)}
\date{CENG 467 — Spring 2026 \\
      İzmir Institute of Technology \\
      Prof. Dr.\ Aytuğ Onan \\[0.5em]
      \today}

\begin{document}
\maketitle
\vspace{-1em}

\section*{1.\ Project Overview}

This project investigates \textbf{controllable text generation with continuous
diffusion language models} (Diffusion-LM, Li et al.\ 2022) and compares the
proposed approach with two autoregressive (AR) baselines on the
\textbf{E2E NLG cleaned} benchmark. The motivating question is whether a
non-autoregressive diffusion-based generator can match the fluency of AR
decoders while delivering measurably higher \emph{lexical diversity} and
\emph{attribute controllability}.

\paragraph{Task.} Given a meaning representation (MR) of restaurant
attributes (e.g.\ \texttt{name[Loch Fyne], food[Japanese], priceRange[high]}),
generate a fluent natural-language description that faithfully realizes all
slots.

\paragraph{Why diffusion?} AR sampling proceeds token-by-token and tends to
collapse onto a narrow part of the output distribution. Diffusion-LM operates
in a continuous embedding space and admits non-autoregressive
\emph{plug-and-play} control via inpainting, classifier guidance, or — in our
implementation — \textbf{prefix clamping}: the tokenized MR occupies a fixed
prefix that is held constant across every denoising step.

\section*{2.\ What is Done (Status by Component)}

\renewcommand{\arraystretch}{1.2}
\begin{tabular}{p{4.6cm} p{1.5cm} p{8.3cm}}
\toprule
\textbf{Component} & \textbf{Status} & \textbf{Notes} \\
\midrule
Project scaffolding         & \textbf{Done} & Public GitHub repo with full directory layout. \\
Data pipeline (E2E NLG)     & \textbf{Done} & Direct CSV load from \texttt{tuetschek/e2e-cleaning}; cached locally. \\
GPT-2 baseline (124M)       & \textbf{Done} & Causal LM fine-tuning, \texttt{<|sep|>}-formatted prompts. \\
T5-small baseline           & \textbf{Done} & Sequence-to-sequence with HF \texttt{Seq2SeqTrainer}. \\
Diffusion-LM model          & \textbf{Done} & 6-layer denoiser, sqrt schedule, x0-prediction, prefix conditioning. \\
\quad — three noise schedules & \textbf{Done} & sqrt / cosine / linear (configurable). \\
\quad — DDIM sampling         & \textbf{Done} & Prefix embeddings clamped at every step. \\
Evaluation                  & \textbf{Done} & BLEU, ROUGE-L, distinct-1/2, slot accuracy. \\
Ablation suite              & \textbf{Done} & 5 configs (schedule, DDIM steps, prefix length). \\
Smoke test                  & \textbf{Done} & 10-second CPU sanity test; runs in CI / Colab. \\
LNCS report skeleton        & \textbf{Done} & All 8 required sections drafted with project content. \\
Error analysis notebook     & \textbf{Done} & Per-slot failure rates, length distribution, qualitative dumps. \\
\midrule
Full training runs          & \textbf{Pending} & Running on Colab Pro (A100). \\
Final numbers in tables     & \textbf{Pending} & Auto-aggregated via \texttt{src/compare\_results.py}. \\
Final paper polish          & \textbf{Pending} & After numbers are filled in. \\
Live demo                   & \textbf{Pending} & Notebook-based; will reuse Colab pipeline. \\
\bottomrule
\end{tabular}

\section*{3.\ Implementation Highlights}

\subsection*{3.1\ Prefix-conditioned Diffusion-LM}

We embed tokens $w_{1:L}$ with a learned embedding matrix and add Gaussian
noise in embedding space:
\[
  q(x_t \mid x_0) = \mathcal{N}\!\big(\sqrt{\bar\alpha_t}\, x_0,\; (1-\bar\alpha_t)\,I\big).
\]
The first $P$ positions hold the tokenized MR and are \emph{never noised}
during training and \emph{never overwritten} during sampling. The denoiser
$f_\theta(x_t, t)$ predicts $x_0$, and the loss combines an embedding-space
MSE with an auxiliary rounding cross-entropy, both restricted to the target
positions:
\[
  \mathcal{L} = \mathbb{E}_{t,\,x_0,\,\epsilon}\!\left[
    \mathbf{1}_\text{tgt}\!\odot\!\big(\|x_0 - f_\theta(x_t,t)\|^2
    + \mathrm{CE}(W^{\!\top}\!f_\theta(x_t,t),\,w)\big)
  \right].
\]
This gives ``hard'' controllability: the model cannot escape the MR.

\subsection*{3.2\ Baselines}

\begin{itemize}[leftmargin=1.2em,topsep=2pt,itemsep=0pt]
  \item \textbf{GPT-2 (124M).} Causal LM with input \texttt{"MR <|sep|> reference\textless{}eos\textgreater{}"}. Decoded with nucleus sampling ($p\!=\!0.95$, $\tau\!=\!1.0$).
  \item \textbf{T5-small.} Encoder--decoder with source = MR text, target = reference. Decoded with beam search ($b\!=\!4$).
\end{itemize}

\subsection*{3.3\ Reproducibility}

The full pipeline runs end-to-end via a single Colab notebook
(\texttt{notebooks/colab\_train.ipynb}) with a \textbf{10-second smoke test}
that exercises the MR parser, all noise schedules, Diffusion-LM forward /
backward / sampling, and the metric pipeline. A separate \texttt{HANDOFF.md}
document allows a collaborator to reproduce the entire run on Colab Pro
without further guidance.

\section*{4.\ Evaluation Plan}

\renewcommand{\arraystretch}{1.15}
\begin{tabular}{p{2.6cm} p{4cm} p{8.5cm}}
\toprule
\textbf{Axis} & \textbf{Metrics} & \textbf{Purpose} \\
\midrule
Fluency           & BLEU, ROUGE-L              & Surface-level fidelity to references. \\
Diversity         & distinct-1, distinct-2     & Lexical variety across the test set. \\
Controllability   & Slot accuracy              & Fraction of MR slot values surfaced verbatim in the output. \\
Sensitivity       & 5-way ablation             & Noise schedule, DDIM steps, prefix length. \\
\bottomrule
\end{tabular}

\section*{5.\ Preliminary Results}

Full training is in progress on Colab Pro (A100). Numbers will be inserted into
the placeholder tables below once the run completes; the same numbers feed
directly into the LNCS final paper via \texttt{results/summary.md}.

\begin{center}
\begin{tabular}{lccccc}
\toprule
Model & BLEU & ROUGE-L & dist-1 & dist-2 & Slot Acc. \\
\midrule
GPT-2 (nucleus)         & — & — & — & — & — \\
T5 (beam-4)             & — & — & — & — & — \\
Diffusion-LM (proposed) & — & — & — & — & — \\
\bottomrule
\end{tabular}
\end{center}

\section*{6.\ Risks and Mitigations}

\begin{itemize}[leftmargin=1.2em,topsep=2pt,itemsep=2pt]
  \item \textbf{Long training runtime.} Mitigation: per-epoch checkpointing,
        smoke test before kickoff, partial training fallback that retains
        baselines only.
  \item \textbf{Diffusion-LM may underperform AR baselines on raw BLEU.}
        Expected per the literature; the project explicitly studies the
        \emph{diversity / control vs.\ fluency} tradeoff, so a small BLEU gap
        is acceptable if compensated by distinct-$n$ and slot-accuracy gains.
  \item \textbf{Slot accuracy is currently exact substring match.} Mitigation:
        for the final paper we will optionally integrate the official E2E
        scorer (BLEU, NIST, METEOR, ROUGE-L, CIDEr) from
        \texttt{tuetschek/e2e-metrics} for stronger external validity.
\end{itemize}

\section*{7.\ Remaining Timeline}

\renewcommand{\arraystretch}{1.15}
\begin{tabular}{p{2.5cm} p{12.5cm}}
\toprule
\textbf{Week} & \textbf{Deliverable} \\
\midrule
12 (now) & Submit this progress report. Complete full Colab training run
            (GPT-2 + T5 + Diffusion-LM). \\
13       & Run the 5 ablation experiments; populate the results tables in
            \texttt{report/main.tex}. \\
14       & Error analysis (per-slot failure modes, qualitative case study);
            iterate on the LNCS draft. \\
15       & Final paper submission and live demo (Colab walkthrough). \\
\bottomrule
\end{tabular}

\section*{8.\ Repository}

The complete project (source, configs, notebooks, report drafts, handoff
documentation) is publicly available at:

\begin{center}
\url{https://github.com/zubeyralmaho/diffusion-lm-ctg}
\end{center}

The repository follows a clear, reproducible structure
(\texttt{src/}, \texttt{configs/}, \texttt{scripts/}, \texttt{notebooks/},
\texttt{report/}, \texttt{tests/}) and includes a \texttt{requirements.txt}
plus end-to-end instructions in \texttt{README.md} and \texttt{HANDOFF.md}.

\end{document}
